\pagestyle{mystyle}
% =============================================================
% =============================================================
% =============================================================
\chapter{Fondamenti Tecnologici del Progetto RAG Aziendale}
\label{chap:background}

Il sistema sviluppato durante il progetto aziendale nasce da un'esigenza
applicativa precisa: rendere interrogabile la conoscenza prodotta in i3Guild
senza trasferire dati interni verso servizi esterni. Prima di descrivere il caso
d'uso, i requisiti e l'implementazione, è quindi necessario introdurre i
concetti tecnici minimi su cui si fonda la soluzione: modelli linguistici
generativi, inferenza locale, \textit{Retrieval-Augmented Generation} (RAG),
embedding, database vettoriali e orchestrazione della pipeline.

Il capitolo non intende esaurire lo stato dell'arte sull'inferenza locale o
sulla quantizzazione, temi che saranno ripresi in modo specifico nei
Capitoli~\ref{chap:ottimizzazione} e~\ref{chap:evaluation}. In questa sede
vengono introdotti solo gli elementi necessari a leggere il prototipo
realizzato in azienda e a comprendere le decisioni progettuali discusse nei
capitoli successivi.

% -------------------------------------------------------------
\section{Architetture Transformer}
\label{sec:transformer}
% -------------------------------------------------------------

\subsection{Dalle reti neurali ricorrenti ai Transformer}
Il modello Transformer, introdotto da Vaswani et al.\ nel
2017~\cite{vaswani2017attention}, costituisce il punto di svolta da cui derivano
gli attuali modelli linguistici. Prima della sua introduzione, gran parte dei
sistemi di trattamento del linguaggio naturale (\textit{Natural Language
Processing}, NLP) si basava su reti neurali ricorrenti (\textit{Recurrent
Neural Networks}, RNN) e, in particolare, su architetture LSTM (\textit{Long
Short-Term Memory})~\cite{hochreiter1997lstm}. Tali modelli elaboravano la
sequenza un elemento alla volta, aggiornando progressivamente uno stato nascosto.
Questa impostazione era coerente con la natura ordinata del testo, ma introduceva
due limiti rilevanti: la scarsa parallelizzabilità dell'addestramento e la
difficoltà nel preservare dipendenze informative tra elementi molto distanti
nella sequenza (\textit{vanishing gradient}).

Il Transformer affronta tali limiti sostituendo la ricorrenza con un meccanismo
di \textit{self-attention} globale, che mette ogni token in relazione diretta
con tutti gli altri token della sequenza, indipendentemente dalla loro distanza
posizionale. Da questa impostazione derivano sia la maggiore efficienza in
addestramento sia la capacità dei modelli successivi di costruire
rappresentazioni contestuali più ricche.

\subsection{Il meccanismo di Self-Attention}
\label{subsec:attention}

Il nucleo computazionale del Transformer è il meccanismo di \textit{attention}
scalato. Data una sequenza di token rappresentati come vettori di dimensione
$d_{\text{model}}$, si calcolano tre proiezioni lineari per ogni elemento:
la query $\mathbf{Q}$, la chiave $\mathbf{K}$ e il valore $\mathbf{V}$. Il
punteggio di attenzione~\cite{vaswani2017attention} è definito come:

\begin{equation}
  \text{Attention}(\mathbf{Q}, \mathbf{K}, \mathbf{V})
  = \text{softmax}\!\left(\frac{\mathbf{Q}\mathbf{K}^\top}{\sqrt{d_k}}\right)\mathbf{V}
  \label{eq:attention}
\end{equation}

dove $d_k$ è la dimensione dei vettori di chiave; il fattore $1/\sqrt{d_k}$
previene la saturazione della funzione softmax in presenza di prodotti scalari
molto grandi.

Il medesimo calcolo viene poi replicato su più \textit{teste} di attenzione
(\textit{attention heads}). Ciascuna testa apprende una proiezione diversa dello
spazio di input e può quindi specializzarsi su relazioni differenti, ad esempio
sintattiche, semantiche o posizionali~\cite{vaswani2017attention}. La
concatenazione degli output di $h$ teste definisce il meccanismo di
\textit{Multi-Head Attention}:

\begin{equation}
  \text{MultiHead}(\mathbf{Q}, \mathbf{K}, \mathbf{V})
  = \text{Concat}(\text{head}_1, \ldots, \text{head}_h)\,\mathbf{W}^O
  \label{eq:multihead}
\end{equation}

con $\text{head}_i = \text{Attention}(\mathbf{Q}\mathbf{W}^Q_i,
\mathbf{K}\mathbf{W}^K_i, \mathbf{V}\mathbf{W}^V_i)$.

La posizione dei token nella sequenza non è catturata dall'attention; viene
iniettata tramite \textit{positional encoding}, ossia vettori sommati agli embedding
di input che forniscono informazioni sull'ordine relativo degli elementi.

\subsection{Architettura encoder-decoder e varianti decoder-only}

L'architettura originale del Transformer è di tipo encoder-decoder. L'encoder
trasforma la sequenza di input in una rappresentazione contestuale densa; il
decoder utilizza tale rappresentazione, tramite \textit{cross-attention}, per
produrre la sequenza di output. Questa struttura è particolarmente adatta a
compiti \textit{sequence-to-sequence}, come la traduzione
automatica~\cite{vaswani2017attention}.

I modelli linguistici generativi contemporanei adottano prevalentemente un'architettura \textit{decoder-only}, nella quale l'attenzione è applicata in modo
\textit{causale} (o \textit{masked}): ogni token dipende esclusivamente dai
precedenti. Questo vincolo è necessario per l'addestramento autoregressivo,
nel quale il modello impara a predire il token successivo dato il contesto.

\subsection{Limiti in contesto zero-shot e motivazione del RAG}
\label{subsec:llm_limits}

La capacità dei Transformer di modellare il linguaggio non elimina tuttavia un
problema centrale per l'uso aziendale: un LLM pre-addestrato non conosce, per
definizione, i dati interni prodotti dopo l'addestramento e non dispone di un
meccanismo nativo per verificarli. In un dominio come quello considerato in
questa tesi, dove le risposte devono riferirsi a documenti aziendali specifici,
questa caratteristica diventa un limite operativo. La
Tabella~\ref{tab:llm_limits} riassume le criticità principali e le relative
strategie di mitigazione.

\begin{table}[ht]
  \centering
  \caption{Limiti strutturali degli LLM in contesti aziendali e relative strategie di mitigazione.}
  \label{tab:llm_limits}
  \small
  \begin{tabular}{p{3cm}p{5.5cm}p{5.5cm}}
    \toprule
    \textbf{Limitazione} & \textbf{Descrizione} & \textbf{Strategia di mitigazione} \\
    \midrule
    Conoscenza statica
      & I parametri codificano la conoscenza fino al \textit{knowledge cutoff};
        nessun accesso a dati interni o aggiornati dopo l'addestramento.
      & RAG: recupero dinamico da knowledge base esterna, aggiornabile
        senza riaddestrare il modello. \\
    \addlinespace
    Allucinazioni~\cite{ji2023survey}
      & Il modello può generare testo plausibile ma fattualmente errato su informazioni
        assenti dalla distribuzione di addestramento.
      & RAG con \textit{grounding}: la risposta è ancorata a documenti
        recuperati e verificabili. \\
    \addlinespace
    Privacy dei dati
      & Il fine-tuning su modelli cloud può esporre dati riservati a infrastrutture
        di terze parti, con rischi di conformità normativa.
      & Esecuzione on-premise con modelli open-weight: nessuna trasmissione
        di dati verso servizi esterni. \\
    \bottomrule
  \end{tabular}
\end{table}

Da questa analisi deriva la scelta di privilegiare il paradigma RAG, discusso
nella Sezione~\ref{sec:rag}. Rispetto al fine-tuning, RAG permette di aggiornare
la conoscenza disponibile intervenendo sul corpus documentale e non sui pesi del
modello; per questo risulta più adatto a basi di conoscenza aziendali soggette a
modifiche frequenti.

% -------------------------------------------------------------
\section{LLM: panorama e analisi comparativa}
\label{sec:llm}
% -------------------------------------------------------------

Una volta chiariti i limiti dei modelli generativi, occorre distinguere tra le
famiglie di LLM disponibili. A partire dal 2022 l'ecosistema si è diviso lungo
una linea rilevante per questa tesi: da un lato i modelli proprietari, accessibili
prevalentemente tramite API cloud; dall'altro i modelli open-weight, i cui pesi
possono essere scaricati ed eseguiti su infrastruttura locale, nel rispetto delle
relative licenze. La Tabella~\ref{tab:llm_comparison} offre una panoramica dei
modelli considerati come riferimento tecnologico.

\begin{table}[ht]
  \centering
  \caption{Panorama comparativo dei principali LLM (2024--2025). \textsuperscript{†}~Esecuzione
    garantita on-premise senza dipendenze cloud; \textsuperscript{‡}~pesi
    scaricabili pubblicamente; MoE~=~\textit{Mixture of Experts}.}
  \label{tab:llm_comparison}
  \small
  \begin{tabular}{llllp{2cm}cc}
    \toprule
    \textbf{Modello} & \textbf{Sviluppatore} & \textbf{Licenza}
      & \textbf{Accesso} & \textbf{Architettura / Param. attivi}
      & \textbf{Ctx} & \textbf{Multi-modale} \\
    \midrule
    GPT-5 / o3       & OpenAI        & Proprietaria  & API cloud
      & Dense, n.d.          & 400k  & $\checkmark$ \\
    Claude 4 Sonnet  & Anthropic     & Proprietaria  & API cloud
      & Dense, n.d.          & 200k  & $\checkmark$ \\
    Gemini 3 Pro     & Google DM     & Proprietaria  & API cloud
      & Dense, n.d.          & 1M    & $\checkmark$ \\
    \midrule
    Llama 4 Scout    & Meta AI       & Open-weight\textsuperscript{‡}
      & Self-hosted\textsuperscript{†} & MoE, 17B & 10M & $\checkmark$ \\
    Mistral Large 3  & Mistral AI    & Apache 2.0\textsuperscript{‡}
      & Self-hosted\textsuperscript{†} & MoE, 41B & 256k & $\checkmark$ \\
    Qwen3-8B         & Alibaba Cloud & Apache 2.0\textsuperscript{‡}
      & Self-hosted\textsuperscript{†} & Dense, 8B & 128k & -- \\
    \bottomrule
  \end{tabular}
\end{table}

\subsection{Modelli proprietari}

I modelli proprietari offrono prestazioni elevate sui benchmark generali, ma
presentano vincoli rilevanti per l'adozione in contesti con requisiti di
riservatezza.

\paragraph{GPT-5 e OpenAI o-series.}
La famiglia GPT-5~\cite{openai2025gpt5} rappresenta un riferimento recente per
pianificazione complessa, sviluppo software e integrazione multimodale. La
serie "o" (\texttt{o1}, \texttt{o3}) è invece orientata a compiti che
richiedono ragionamento deliberato, con miglioramenti osservati in ambiti
scientifici e logico-matematici~\cite{openai2024o1}. L'erogazione è vincolata
ad API cloud proprietarie; in ambito enterprise è possibile operare tramite
servizi dedicati, ad esempio Azure OpenAI Service, che offrono garanzie
contrattuali di isolamento dei dati.

\paragraph{Claude 4 (Anthropic).}
I modelli Claude adottano approcci come la \textit{Constitutional AI}~\cite{bai2022constitutional}, un protocollo di addestramento guidato da principi
per migliorare l'allineamento. Anche questi modelli sono erogati in modalità
cloud; Anthropic propone modalità contrattuali (ad esempio \textit{Zero Data Retention})
per rispondere a esigenze enterprise~\cite{anthropic2025claude4}.

\paragraph{Gemini (Google DeepMind).}
La famiglia Gemini~\cite{google2025gemini3} è caratterizzata da multimodalità
nativa e da finestre di contesto molto estese, fino a un milione di token per
Gemini 3 Pro.
Anche in questo caso l'accesso è fornito tramite servizi cloud.

\subsection{Motivazioni per l'adozione di modelli locali}
\label{subsec:onpremise_motivations}

Nel caso in esame, le prestazioni assolute non sono l'unico criterio di scelta.
Il sistema descritto in questa tesi deve poter operare su dati aziendali interni
e deve quindi privilegiare controllo, tracciabilità e assenza di dipendenze da
servizi esterni. Per questo l'architettura utilizza modelli eseguiti interamente
in locale (\textit{self-hosted}). La scelta si fonda su tre motivazioni,
sintetizzate nella Tabella~\ref{tab:onpremise_motivations}.

\begin{table}[ht]
  \centering
  \caption{Motivazioni per l'adozione di modelli LLM on-premise.}
  \label{tab:onpremise_motivations}
  \small
  \begin{tabular}{p{3.5cm}p{5.5cm}p{4.5cm}}
    \toprule
    \textbf{Motivazione} & \textbf{Descrizione} & \textbf{Riferimento normativo / letteratura} \\
    \midrule
    Sovranità del dato e conformità
      & L'esecuzione locale evita la trasmissione di dati verso infrastrutture
        di terze parti, anche in presenza di garanzie contrattuali (ZDR, DPA).
      & GDPR~\cite{regulation2016gdpr}; rischi del cloud nella gestione di dati sensibili. \\
    \addlinespace
    Costi operativi nel lungo periodo
      & Il pricing a consumo dei provider cloud può risultare più oneroso rispetto
        all'ammortamento dell'hardware locale in scenari ad alto volume.
      &~\cite{bommasani2021opportunities}; rischio di \textit{vendor lock-in}. \\
    \addlinespace
    Controllo e personalizzazione
      & I modelli open-weight (LLaMA, Mistral, Qwen) consentono fine-tuning
        supervisionato e allineamento su dati propri, senza le restrizioni delle policy
        dei sistemi proprietari.
      &~\cite{touvron2023llama,jiang2023mistral,qwen2024} \\
    \bottomrule
  \end{tabular}
\end{table}

Per queste ragioni, i modelli proprietari sono mantenuti come riferimento
comparativo, mentre l'analisi progettuale si concentra sulle soluzioni
eseguibili localmente. Tale impostazione rende coerente la scelta del modello
con i vincoli di privacy, costo e governo dell'infrastruttura discussi nei
capitoli successivi.

\subsection{Modelli open-weight per uso on-premise}

\paragraph{Llama 4 (Meta AI).}
Rilasciata nell'aprile 2025, la famiglia Llama 4~\cite{meta2025llama4} adotta
architetture \textit{Mixture of Experts} (MoE) e supporta la multimodalità.
Le varianti \textit{Scout} e \textit{Maverick} bilanciano parametri effettivi e
numero di esperti per gestire contesti estesi fino a 10 milioni di token. La
licenza open-weight presenta clausole specifiche da valutare in fase di adozione,
in particolare per operatori e per l'Unione Europea.

\paragraph{Mistral 3 (Mistral AI).}
Nel 2025 Mistral AI ha ampliato la propria offerta con modelli orientati
all'integrazione locale~\cite{mistral2025large3}. Il modello di punta,
\textit{Mistral Large 3}, è un modello MoE multilingue; la famiglia include
varianti ottimizzate per l'esecuzione su hardware locale. La distribuzione sotto
licenza Apache 2.0 facilita l'impiego in progetti aperti.

\paragraph{Qwen3 (Alibaba Cloud).}
La generazione Qwen3~\cite{qwen2025qwen3} introduce modalità operative differenziate
per task logici ed efficienza. La famiglia copre modelli da dimensioni molto
ridotte fino a varianti MoE con finestre di contesto estese. Nel Capitolo~\ref{chap:implementazione}
si descrive la scelta della variante 8B per il sistema proposto, motivata dal
bilancio tra qualità semantica e velocità di inferenza su CPU-only.

\subsection{Inferenza locale nel prototipo aziendale}

L'esecuzione locale di un LLM introduce un vincolo ulteriore rispetto all'uso di
API cloud: il modello deve convivere con database, servizi applicativi,
pipeline RAG e runtime di inferenza nello stesso ambiente controllato. Nel
progetto i3Guild questo requisito deriva dalla necessità di non esporre dati
aziendali a servizi esterni. La conseguenza pratica è che la scelta del modello
non può basarsi solo sulla qualità linguistica, ma deve considerare memoria
disponibile, latenza di caricamento, tempo al primo token e throughput.

La quantizzazione viene quindi introdotta qui come motivazione tecnica generale:
ridurre la precisione dei pesi può diminuire l'occupazione di memoria e rendere
eseguibili modelli altrimenti troppo costosi. Tuttavia, l'effetto reale dipende
dal runtime, dal formato del modello e dal costo di dequantizzazione. Per questo
l'analisi sistematica della quantizzazione viene separata dal progetto
aziendale e ripresa nei capitoli finali, dedicati all'inferenza locale su Apple
Silicon.

\subsection{Ollama: server di inferenza locale}

Alla luce dei vincoli descritti, il progetto adotta un'infrastruttura di
inferenza locale gestita tramite
\textbf{Ollama}~\cite{ollama2023}, un server open source che espone modelli
locali in formato GGUF tramite un'API REST compatibile con la specifica
\textit{OpenAI Chat Completions}. Ollama semplifica il ciclo di vita dei modelli
locali: l'installazione avviene tramite un comando semplice (\texttt{ollama pull <modello>}),
la gestione del contesto e del formato di prompt è automatizzata, e il server
può essere distribuito come container Docker. Nel contesto di questo progetto,
Ollama serve il modello di chat; la generazione degli embedding è invece
affidata a una pipeline separata basata su FastEmbed e Haystack, come descritto
nella Sezione~\ref{subsec:embedding_models}.

% -------------------------------------------------------------
\section{Retrieval-Augmented Generation (RAG)}
\label{sec:rag}
% -------------------------------------------------------------

\subsection{Motivazione e definizione}

Il paradigma \textit{Retrieval-Augmented Generation} (RAG), formalizzato da
Lewis et al.~\cite{lewis2020retrieval}, risponde al limite degli LLM come
sistemi a conoscenza chiusa (\textit{closed-book}). In un modello generativo
tradizionale, infatti, la conoscenza disponibile è codificata nei parametri e
non può essere aggiornata senza un nuovo ciclo di addestramento o fine-tuning.

In un sistema RAG, la generazione viene preceduta da una fase di
\textit{retrieval} dinamico. Prima di invocare l'LLM, il sistema interroga una
base di conoscenza esterna e recupera i documenti più pertinenti rispetto alla
query. Tali documenti vengono poi inseriti nel prompt
(\textit{prompt augmentation}), così che la risposta sia fondata non solo sulla
conoscenza parametrica del modello, ma anche su fonti aggiornate e
verificabili~\cite{shuster2021retrieval}. Questa separazione tra modello e
conoscenza esterna è il motivo per cui RAG si adatta bene a domini aziendali in
evoluzione.

\subsection{Architettura del sistema RAG}
\label{subsec:rag_arch}

Dal punto di vista architetturale, un sistema RAG separa il momento in cui la
conoscenza viene preparata dal momento in cui viene interrogata. Ne derivano due
pipeline distinte, messe a confronto nella Tabella~\ref{tab:rag_pipelines} e
illustrate nella Figura~\ref{fig:rag_pipeline}: la pipeline di
\textbf{ingestione}, eseguita offline, e la pipeline di \textbf{inferenza},
attivata in tempo reale a ogni richiesta utente.

\begin{table}[ht]
  \centering
  \caption{Le due pipeline del sistema RAG a confronto.}
  \label{tab:rag_pipelines}
  \small
  \begin{tabular}{p{2.5cm}p{5.5cm}p{5.5cm}}
    \toprule
    \textbf{Fase} & \textbf{Ingestion pipeline (offline)}
      & \textbf{Query pipeline (real-time)} \\
    \midrule
    1 & Pre-processing e chunking dei documenti sorgente
      & Embedding della query utente nello stesso spazio vettoriale dei documenti \\
    \addlinespace
    2 & Generazione di embedding densi e sparsi tramite modelli locali
      & Ricerca ibrida nel database vettoriale (\textit{top-k}) con filtri su metadati \\
    \addlinespace
    3 & Indicizzazione in Qdrant con payload di metadati strutturati
      & Augmented generation: contesto recuperato + query $\rightarrow$ LLM \\
    \addlinespace
    \textit{Trigger}  & Esecuzione manuale o schedulata (es. giornaliera)
      & Ad ogni richiesta utente \\
    \addlinespace
    \textit{Latenza}  & Minuti / ore (elaborazione batch)
      & Millisecondi (inferenza real-time) \\
    \bottomrule
  \end{tabular}
\end{table}

\subsubsection{Pipeline di ingestione}

La pipeline di ingestione trasforma i documenti grezzi in una rappresentazione
adatta alla ricerca semantica. Il risultato non è più soltanto un archivio di
testi, ma un indice interrogabile per similarità. Il documento viene anzitutto
normalizzato e, quando necessario, suddiviso in unità testuali di dimensione
controllata. La granularità del chunking incide direttamente sulla qualità del
retrieval: chunk troppo piccoli perdono contesto, mentre chunk troppo grandi
possono diluire il segnale semantico o superare la finestra del modello di
embedding~\cite{chen2023benchmarking}.

Ogni chunk viene poi trasformato in due rappresentazioni complementari. La
prima è un vettore denso prodotto da un modello di embedding
$f\colon \mathcal{T} \rightarrow \mathbb{R}^d$, in modo che documenti
semanticamente simili producano vettori vicini secondo la similarità coseno:

\begin{equation}
  \text{sim}(\mathbf{u}, \mathbf{v})
  = \frac{\mathbf{u} \cdot \mathbf{v}}{\|\mathbf{u}\| \cdot \|\mathbf{v}\|}
  \label{eq:cosine}
\end{equation}

La rappresentazione sparsa descrive invece il testo attraverso pesi associati a
termini o feature lessicali. Questa componente conserva segnali che l'embedding
denso può attenuare, come nomi propri, tecnologie, sigle o termini specifici del
dominio.

Le rappresentazioni dense e sparse vengono infine persistite in
\textbf{Qdrant}~\cite{qdrant2023}, insieme ai metadati strutturati necessari al
filtraggio. Per la componente densa Qdrant usa l'algoritmo HNSW
(\textit{Hierarchical Navigable Small World})~\cite{malkov2018efficient} per la
ricerca approssimata (\textit{Approximate Nearest Neighbor}, ANN), mentre la
componente sparsa abilita il matching lessicale pesato. La scelta di Qdrant
rispetto ad alternative quali pgvector, Chroma e Milvus è motivata nel
Capitolo~\ref{chap:progettazione}.

\subsubsection{Pipeline di inferenza}

La pipeline di inferenza utilizza l'indice costruito offline per rispondere a
una query utente. A differenza dell'ingestione, questa fase è vincolata dalla
latenza percepita e deve bilanciare qualità del recupero, numero di documenti
restituiti e costo della generazione. La query viene trasformata nelle stesse
due rappresentazioni usate in indicizzazione: un embedding denso per la
similarità semantica e una rappresentazione sparsa per il matching lessicale.
Su queste rappresentazioni si esegue una ricerca ibrida in Qdrant, combinando
il punteggio denso e quello sparso; quando la query contiene vincoli
strutturati, il retrieval può essere limitato tramite \textit{metadata
filtering}, ad esempio ai documenti appartenenti a una specifica epoca
temporale. Il contesto recuperato viene quindi concatenato alla query nel prompt
inviato all'LLM, che genera la risposta usando sia la propria conoscenza
parametrica sia le informazioni esterne fornite dal sistema.

\begin{figure}[ht]
  \centering
  \resizebox{\textwidth}{!}{%
  \begin{tikzpicture}[
  node distance=1.0cm and 1.15cm,
  box/.style={
    draw,
    rounded corners=3pt,
    align=center,
    minimum width=2.8cm,
    minimum height=0.85cm,
    font=\small
  },
  widebox/.style={
    draw,
    rounded corners=3pt,
    align=center,
    minimum width=3.9cm,
    minimum height=0.9cm,
    font=\small
  },
  store/.style={
    draw,
    rounded corners=4pt,
    align=center,
    minimum width=4.1cm,
    minimum height=1.15cm,
    font=\small,
    thick
  },
  arrow/.style={->, thick},
  group/.style={
    draw,
    rounded corners=6pt,
    inner sep=0.35cm,
    dashed
  }
]

% --------------------
% Ingestion pipeline
% --------------------
\node[box] (docs) {Documenti\\sorgente};
\node[box, right=of docs] (prep) {Pulizia e\\chunking};
\node[box, right=of prep] (docdense) {Embedding\\denso};
\node[box, below=of docdense] (docsparse) {Embedding\\sparso};

\node[store, right=1.4cm of docdense, yshift=-0.45cm] (qdrant)
{Qdrant\\
\footnotesize vettori densi, sparsi\\
\footnotesize metadati e payload};

% --------------------
% Inference pipeline
% --------------------
\node[box, below=2.1cm of docs] (query) {Query\\utente};
\node[box, right=of query] (qrewrite) {Normalizzazione\\query};
\node[box, right=of qrewrite] (qdense) {Embedding\\denso};
\node[box, below=of qdense] (qsparse) {Embedding\\sparso};

\node[widebox, right=1.4cm of qdense, yshift=-0.45cm] (retrieval)
{Retrieval ibrido\\
\footnotesize score denso + sparso\\
\footnotesize filtri sui metadati};

\node[widebox, right=of retrieval] (prompt)
{Prompt aumentato\\
\footnotesize contesto recuperato + query};

\node[box, right=of prompt] (llm) {LLM\\locale};
\node[box, right=of llm] (answer) {Risposta\\fondata};

% --------------------
% Arrows: ingestion
% --------------------
\draw[arrow] (docs) -- (prep);
\draw[arrow] (prep) -- (docdense);
\draw[arrow] (prep) |- (docsparse);
\draw[arrow] (docdense) -- (qdrant);
\draw[arrow] (docsparse) -- (qdrant);

% --------------------
% Arrows: inference
% --------------------
\draw[arrow] (query) -- (qrewrite);
\draw[arrow] (qrewrite) -- (qdense);
\draw[arrow] (qrewrite) |- (qsparse);
\draw[arrow] (qdense) -- (retrieval);
\draw[arrow] (qsparse) -- (retrieval);
\draw[arrow] (qdrant) -- (retrieval);
\draw[arrow] (retrieval) -- (prompt);
\draw[arrow] (prompt) -- (llm);
\draw[arrow] (llm) -- (answer);

% --------------------
% Labels
% --------------------
\node[font=\small\bfseries, above=0.15cm of prep]
{Pipeline di ingestione};

\node[font=\small\bfseries, above=0.15cm of qrewrite]
{Pipeline di inferenza};

\end{tikzpicture}
  }
  \caption{Schema della pipeline RAG ibrida: la fase di ingestione costruisce
    rappresentazioni dense, sparse e metadati; la fase di inferenza combina
    retrieval semantico, matching lessicale e filtri strutturati prima della
    generazione aumentata~\cite{lewis2020retrieval}.}
  \label{fig:rag_pipeline}
\end{figure}


\subsection{Modelli di embedding}
\label{subsec:embedding_models}

La qualità del retrieval dipende in larga misura dal modello di embedding: se
la rappresentazione vettoriale non preserva correttamente il significato dei
documenti, anche il modello generativo riceverà un contesto poco pertinente. Il
vincolo di esecuzione on-premise restringe quindi la scelta a modelli
distribuibili localmente. La versione aggiornata del progetto usa FastEmbed
tramite Haystack e produce due segnali complementari: embedding densi, adatti
alla similarità semantica, ed embedding sparsi, utili per il matching di termini
specifici. L'analisi completa della configurazione adottata è presentata nel
Capitolo~\ref{chap:progettazione}.

\subsection{Database vettoriali}
\label{subsec:vector_db}

Un database vettoriale è un sistema di storage specializzato per vettori densi
ad alta dimensionalità, progettato per rispondere a query di tipo
\textit{k-nearest neighbor} (kNN). A differenza dei database relazionali,
dove le query sono esatte, la ricerca vettoriale è per natura approssimata:
si cercano i vettori più vicini accettando un trade-off tra richiamo
(\textit{recall}) e velocità (\textit{throughput}).

L'algoritmo HNSW costruisce un grafo gerarchico multi-livello nello spazio
vettoriale, con complessità di ricerca approssimativamente $O(\log n)$ rispetto
al numero di punti $n$ — contro $O(n)$ della ricerca lineare esatta~\cite{malkov2018efficient} — garantendo scalabilità anche per ampi \textit{corpora}.


\subsection{Agentic RAG}
\label{subsec:agentic_rag}

Il RAG lineare descritto finora segue un flusso deterministico: recupero dei
documenti, costruzione del prompt e generazione della risposta. Nei casi in cui
la query richiede più passaggi logici, confronto tra fonti o riformulazioni
successive, tale schema può risultare insufficiente. L'\textit{Agentic RAG}
estende questo approccio utilizzando l'LLM come unità di ragionamento capace di
invocare strumenti, pianificare ricerche successive e valutare la pertinenza dei
documenti recuperati. Architetture come ReAct~\cite{yao2023react} formalizzano
questo ciclo combinando ragionamento verbale e azioni.

In contesti RAG avanzati, tecniche come Self-RAG~\cite{asai2023selfrag} e
Corrective RAG (CRAG)~\cite{yan2024corrective} introducono meccanismi di
controllo sulla qualità del retrieval e della risposta generata. Il sistema può
così rilevare lacune informative, richiedere nuovi documenti o modificare il
prompt prima di produrre la risposta finale. Le rassegne più recenti
evidenziano che questa evoluzione è particolarmente utile quando le query
richiedono sintesi da fonti eterogenee o passaggi intermedi di
ragionamento~\cite{gao2024retrieval}.

Un'altra direzione di sviluppo è rappresentata da \textbf{GraphRAG},
formalizzato da Microsoft Research nel 2024~\cite{edge2024graphrag}. A
differenza del RAG vettoriale classico, che è efficace nel recupero di passaggi
localmente simili alla query ma meno adatto a interrogazioni globali sull'intero
dataset, GraphRAG estrae entità e relazioni dai testi sorgente e costruisce un
grafo di conoscenza gerarchico. In fase di retrieval, il sistema può quindi
aggregare riassunti semantici di intere \textit{community} di concetti,
ottenendo un livello di astrazione più adeguato all'analisi di corpora
documentali estesi.

\subsection{Tecniche di adattamento al dominio}
\label{subsec:domain_adaptation}

Il collegamento tra LLM e conoscenza aziendale può essere realizzato in modi
diversi. RAG mantiene separati i pesi del modello e il corpus documentale;
il fine-tuning, invece, modifica il modello affinché si adatti a un dominio o a
uno stile di risposta specifico. Questa sezione introduce le principali tecniche
di adattamento a parametri ridotti, utili per inquadrare una scelta progettuale
del lavoro: mantenere il corpus aziendale esterno ai pesi del modello e
aggiornabile tramite pipeline di ingestione.

\subsubsection{Fine-tuning con adattatori: LoRA}

Il fine-tuning completo di un LLM richiede risorse computazionali elevate.
Hu et al.~\cite{hu2022lora} hanno proposto \textit{Low-Rank Adaptation} (LoRA),
che vincola gli aggiornamenti dei pesi a matrici di rango basso. Dato un layer
lineare con matrice $\mathbf{W}_0 \in \mathbb{R}^{m \times n}$, LoRA introduce:

\begin{equation}
  \mathbf{W} = \mathbf{W}_0 + \Delta\mathbf{W} = \mathbf{W}_0 + \mathbf{B}\mathbf{A}
  \label{eq:lora}
\end{equation}

con $\mathbf{A} \in \mathbb{R}^{r \times n}$, $\mathbf{B} \in \mathbb{R}^{m
\times r}$ e $r \ll \min(m,n)$. I pesi originali $\mathbf{W}_0$ rimangono
congelati; solo $\mathbf{A}$ e $\mathbf{B}$ vengono ottimizzati. Con $r = 8$
e $m = n = 4096$, il numero di parametri addestrabili si riduce di circa due
ordini di grandezza rispetto al fine-tuning completo.

\subsubsection{QLoRA: fine-tuning a 4-bit}

\textit{Quantized LoRA} (QLoRA)~\cite{dettmers2023qlora} combina la
quantizzazione a 4-bit del modello base con adattatori LoRA a precisione piena.
I pesi vengono quantizzati in formato NF4 (\textit{NormalFloat4}), ottimizzato
per le distribuzioni tipiche dei pesi neurali; durante il forward pass i pesi
vengono de-quantizzati dinamicamente in BF16, mentre gli adattatori operano in BF16.

La Tabella~\ref{tab:lora_vs_qlora} riassume i principali trade-off tra le due
varianti.

\begin{table}[ht]
  \centering
  \caption{Confronto tra LoRA e QLoRA per un modello base da 7B parametri.}
  \label{tab:lora_vs_qlora}
  \small
  \begin{tabular}{lcc}
    \toprule
    \textbf{Caratteristica} & \textbf{LoRA} & \textbf{QLoRA} \\
    \midrule
    Precisione pesi base             & FP16/BF16   & NF4 (4-bit) \\
    Precisione adattatori            & FP16/BF16   & BF16        \\
    Memoria GPU (modello 7B)         & $\sim$14 GB & $\sim$4 GB  \\
    Overhead de-quantizzazione       & assente     & presente    \\
    Degradazione qualitativa stimata & trascurabile & ${<}1\%$ su benchmark principali \\
    \bottomrule
  \end{tabular}
\end{table}

\subsubsection{DoRA: Weight-Decomposed Low-Rank Adaptation}

Nel 2024 Liu et al.\ hanno introdotto \textbf{DoRA}
(\textit{Weight-Decomposed Low-Rank Adaptation})~\cite{liu2024dora}. Il metodo
parte dall'osservazione che LoRA approssima bene l'aggiornamento direzionale dei
pesi, ma cattura con minore efficacia le variazioni di magnitudine osservate nel
fine-tuning completo. DoRA separa quindi il tensore dei pesi in una componente
di magnitudine e una componente direzionale: la prima viene aggiornata
direttamente, mentre la seconda sfrutta l'approssimazione a rango basso di LoRA.
Il divario prestazionale rispetto al fine-tuning completo si riduce così senza
rinunciare al vantaggio principale degli adattatori, cioè il numero contenuto di
parametri addestrabili.

\subsubsection{Prompt engineering e inference-time intervention}

In alternativa al fine-tuning parametrico, è possibile adattare il comportamento
di un LLM intervenendo sulla formulazione dell'input senza modificare i pesi.
Tecniche come \textit{few-shot prompting}~\cite{brown2020language} e
\textit{chain-of-thought prompting}~\cite{wei2022chain} guidano il modello verso
risposte più strutturate, arricchendo il prompt con esempi o istruzioni di
ragionamento. Nel sistema descritto in questa tesi, il prompt engineering è
impiegato nella definizione del \textit{system prompt} e nell'iniezione
strutturata del contesto RAG. Le strategie adottate sono descritte nel
Capitolo~\ref{chap:implementazione}, insieme agli endpoint che costruiscono il
prompt aumentato.

% -------------------------------------------------------------
\section{Framework di orchestrazione: Haystack}
\label{sec:haystack}
% -------------------------------------------------------------

La costruzione di un sistema RAG non consiste soltanto nella scelta del modello
generativo. È necessario coordinare modelli di embedding, database vettoriali,
LLM e fasi di pre-processing attraverso interfacce coerenti e componibili.
\textbf{Haystack}~\cite{haystack2023}, sviluppato da deepset, è un framework
Python open source progettato per organizzare questi componenti in pipeline
esplicite.

Nella versione 2.x adottata in questo progetto, ogni pipeline è un grafo
diretto aciclico (DAG) di componenti con interfacce tipizzate, collegabili in
modo dichiarativo. Sono disponibili integrazioni native con Qdrant
(\texttt{QdrantDocumentStore}), FastEmbed
(\texttt{FastembedDocumentEmbedder}, \texttt{FastembedTextEmbedder} e
componenti sparse corrispondenti), Ollama (\texttt{OllamaGenerator}) e le
principali operazioni di pre-processing (\texttt{DocumentCleaner},
\texttt{DocumentSplitter}).
Il dettaglio implementativo della pipeline di ingestione e di retrieval è
descritto nel Capitolo~\ref{chap:implementazione}.

% -------------------------------------------------------------
\section{Sintesi e posizionamento del lavoro}
\label{sec:background_summary}
% -------------------------------------------------------------

Il capitolo ha definito il quadro tecnico necessario per comprendere il progetto
aziendale descritto nel seguito della tesi. Le architetture Transformer sono il
substrato comune della generazione e degli embedding; i loro limiti nei contesti
a conoscenza chiusa, in particolare conoscenza statica, allucinazioni e assenza
di accesso nativo ai dati interni, spiegano l'adozione del paradigma RAG. La
scelta di operare con modelli e infrastrutture on-premise, come Ollama, Qdrant
e Haystack, risponde invece ai vincoli di riservatezza e manutenibilità del caso
i3Guild. Su queste basi, il capitolo successivo traduce i fondamenti teorici in
requisiti, scelte tecnologiche e architettura del sistema.
